雅化不适合用单一季度 PE 或单一券商目标价估值。锂盐业务具有强周期性，民爆业务又缺少可审计的分部净利和现金流，故主方法是以自有熊/基/牛 EPS 乘周期校准 PE 得到基本面价值，再以可观察市场价格作有限情绪锚；外部卖方目标仅作零权重的敏感性参考。

\section{基本面锚：周期校准 PE，而非 H1 年化}

情景价值为 10.3/17.3/26.1 元，对应熊/基/牛 EPS 1.30/1.74/2.14 元和约 7.9/9.9/12.2 倍 PE。基准倍数低于三家锂企 2026E 平均 PE 12.79 倍，反映雅化尚未证明现金转换及资源成本优势；牛情景才允许部分折价收窄。民爆历史利润池已在 EPS 中体现，但因缺分部净利、项目回款和独立现金流，未另行赋予 SOTP 溢价。

\begin{exhibitbox}[表 8-1：最终估值表]
\begin{tabularx}{\exhibitboxwidth}{L{3.1cm}R{2.2cm}R{2.2cm}R{2.2cm}X}
\toprule
项目 & 熊 & 基准 & 牛 & 注释 \\
\midrule
2026E 归母净利 & 15.0 亿元 & 20.1 亿元 & 24.7 亿元 & 锂盐与民爆的合并情景模型 \\
2026E EPS & 1.30 元 & 1.74 元 & 2.14 元 & 总股本 11.526 亿股 \\
周期校准 PE & 7.9 倍 & 9.9 倍 & 12.2 倍 & 与现金、周期和证据质量匹配 \\
基本面价值 & 10.3 元 & 17.3 元 & 26.1 元 & 不是市场价格预测 \\
概率 & 30\% & 50\% & 20\% & 概率加权内在价值 16.95 元 \\
相对 17.75 元回报 & -42.2\% & -2.4\% & +46.9\% & 截至 2026-07-23 \\
\bottomrule
\end{tabularx}
\end{exhibitbox}

\section{多锚目标：市场价格有信息，但不拥有否决权}

我们以 85\% 权重使用概率加权内在价值 16.95 元，以 15\% 权重使用现价 17.75 元作为可观察市场情绪锚，得到未取整最终目标 17.07 元，展示为 17.0 元。东吴 42 元目标价因只有一份原始报告、且对应更激进的 2026E 归母与 16 倍 PE，保留为零权重外部敏感性；国信和开源没有可用目标价，不能被补成 Street 目标。

\begin{exhibitbox}[表 8-2：基本面、市场与外部锚的权重]
\begin{tabularx}{\exhibitboxwidth}{L{3.4cm}R{2.4cm}R{2.1cm}X}
\toprule
估值锚 & 数值 & 权重 & 使用原因 \\
\midrule
概率加权内在价值 & 16.95 元 & 85\% & 自有情景、EPS 和周期 PE 可复算 \\
市场情绪锚 & 17.75 元 & 15\% & 反映可观察的交易定价与情绪溢价 \\
外部卖方锚 & 东吴 42 元 & 0\% & 单一外部目标；盈利与倍数假设更激进 \\
未取整最终目标 & 17.07 元 & 100\% & 展示为 17.0 元 \\
\bottomrule
\end{tabularx}
\end{exhibitbox}

\section{市场预期桥：现价并不便宜到足以忽略现金风险}

17.75 元分别对应自有熊/基/牛 2026E PE 13.64/10.18/8.28 倍；对应三家外部 EPS 则为 6.75--8.83 倍。低于同行平均倍数可以反映潜在上行，也可以反映现金流、资源成本和 H2 兑现的不确定性。当前价已经接近基准价值，因此目标价上行主要取决于牛情景概率上升，而不是估值倍数的机械扩张。

\begin{sourcequalitybox}[估值失效条件]
若正式 H1 显示现金转换、应收/存货和单位经济支持基准甚至牛情景，当前目标应上修；若利润未能转为现金或 H2 的量价成本路径被削弱，熊情景的概率应提高。报告不把 17.0 元当成精确交易点，而把它作为“现价缺乏安全边际”的概率结论。
\end{sourcequalitybox}
